% ============================================================
% Changes from v2:
%   (f) Co-emergence paragraph split into two:
%       (f1) signature selection (SR level) -- unchanged from v2
%       (f2) NEW: metric rigidity and the GR space/time budget
%            compensation, traced to the Einstein equations
%            emerging from delta S_Pi / delta g_mu_nu = 0.
%            Forward ref to Gravity companion paper.
%       This addresses the question: why does dilating time
%       exactly compensate spatial contraction (ds^2 invariant)?
%       Answer: the signature guarantees existence of a pseudo-
%       Riemannian metric, but it is the spectral entropy
%       functional S_Pi[g] that dynamically locks the metric
%       coefficients through the Einstein field equations,
%       making the compensation exact and position-dependent.
%
% CROSS-REFERENCE UPDATES IN REST OF DOCUMENT:
%   - Section 2 (static chi): forward ref to this section.
%   - Section 7.5 closing sentence: point to Sec. 7.6.
%   - Minkowski/signature section: back-ref to Sec. 7.6.
%   - Gravity companion paper section: note that the dynamical
%     locking of the metric is further discussed in Sec. 7.6.
% ============================================================

\subsection{The Self-Referential Origin of Temporal Succession}
  \label{sec:temporal-succession}

  The framework introduced in the preceding sections places the substrate field~$\chi$
  at a pre-geometric level: no spacetime manifold, no metric, and no notion of physical time is presupposed.
  Observable physics---including time itself---emerges through the admissibility constraints and the projection~$\Pi$.
  This raises a structural question that must be addressed explicitly:
  if $\chi$ is timeless, in what sense does the universe ``evolve,'' and what
  determines the succession from one projected state to the next?
  The present section provides a logically closed account of temporal emergence
  within the framework and clarifies how the spatial and temporal sectors, despite
  having distinct logical origins, remain coupled in the manner required by special
  relativity.

  \paragraph{Two levels: structural rules and their execution.}
    The relational rules encoded by~$\chi$---the Laplacian weights~$w_{ij}$,
    the saturation bound~$c_\chi$, the admissibility envelopes, the fiber
    topology of~$\Pi$---constitute a fixed structural layer.
    They are written once and do not themselves change.
    This layer is analogous to the source code of a deterministic program:
    all rules are simultaneously present, yet the rules alone do not constitute a running computation.
    Physical time belongs to the execution layer, not to the code layer.
    The apparent paradox of a ``timeless'' substrate generating temporal order
    dissolves once these two levels are kept distinct.

  \paragraph{The ordering parameter $\tau$ and the relational energy.}
    The relaxation cascade introduced in Section~\ref{sec:relax-cascade}
    is governed by the flow equation
    \begin{equation}
      \frac{d\varphi_i}{d\tau} = -(\mathcal{L}\,\varphi)_i,
    \end{equation}
    where $\tau$ is a monotonically increasing parameter indexing the sequence
    of configurations.
    No physical interpretation of~$\tau$ as a fundamental time is imposed at this level.
    Instead, $\tau$ labels a strict ordering on the space of relational configurations
    defined by the monotone decrease of the relational energy
    \begin{equation}
      E_\chi[\varphi] \;=\; \|\mathcal{L}\,\varphi\|^2
      \;=\; \sum_{i,j} w_{ij}\,(\varphi_i - \varphi_j)^2,
    \end{equation}
    already defined in Eq.~\eqref{eq:relational-energy}.
    This ordering is intrinsic: it follows from the structure of~$\mathcal{L}$ alone,
    without reference to any background clock.
    Physical time is the coarse-grained, projected image of this ordering,
    once the hyperbolic signature of the effective operator is established
    (\cite[Sec.~6]{Beau2026c}).

  \paragraph{The universe state $U$ as its own cursor.}
    The central self-referential feature of the framework is the following.
    The current macroscopic state of the universe~$U_n$---understood as the totality
    of projected observable data at a given level of the cascade---does not merely
    record the outcome of the preceding transition.
    It determines which successor states~$U_{n+1}$ are admissible.

    More precisely:
    \begin{itemize}
      \item $U_n$ encodes the current spectral structure of the effective Laplacian,
      including the local connectivity weights, the saturation regime, and the
      descriptive resolution scale~$\ell(U_n)$~(cf.~\cite{Beau2026d}).
      \item These data define the admissibility constraints on the next projection:
      which configurations of~$\chi$ project injectively onto observable fields,
      and which local transitions remain within the bounded-flux regime.
      \item Among these constraints, the bounded-flux condition~$c_\chi$ imposes that
      the local relational flux must remain at saturation throughout an admissible
      transition: any descent along~$E_\chi$ that is shallower than the steepest
      admissible gradient would leave residual flux below the saturation bound,
      meaning that the projection retains further resolvable structure and the
      transition is not yet complete.
      Only the steepest admissible descent exhausts the available flux budget and
      closes the transition.
      No external selection principle is involved; the saturation condition itself
      determines the unique realised step.
    \end{itemize}
    There is therefore no external cursor advancing the universe from~$U_n$
    to~$U_{n+1}$.
    The state~$U_n$ is the cursor.
    It carries within itself the conditions that make the next step both determinate
    and admissible.

  \paragraph{The cursor acts on fibre classes, not on $\chi$ directly.}
    Because the projection~$\Pi : \Omega \to \mathcal{O}$ is non-injective, a single
    observable state~$U_n$ corresponds to an entire fibre of substrate
    configurations,
    \begin{equation}
      \Pi^{-1}(U_n) \;\subset\; \Omega.
    \end{equation}
    The cursor does not select a unique microscopic configuration of~$\chi$; it selects
    an admissible fibre class.
    The residual multiplicity within~$\Pi^{-1}(U_n)$ is precisely the source of
    projective fluctuations: the effective observables at level~$U_n$ cannot
    distinguish between the microscopic configurations in the fibre, and this
    irreducible indeterminacy propagates into the next transition.
    Quantum fluctuations, in this picture, are not imposed as an additional
    postulate: they reflect the bounded resolvability of the projection at each step
    of the cascade.

  \paragraph{Irreversibility and the arrow of time.}
    The succession $U_n \to U_{n+1}$ is irreversible by construction:
    $E_\chi$ is non-increasing along admissible transitions, with a strict decrease
    whenever~$U_n$ is not already a stationary configuration of~$\mathcal{L}$.
    This monotone decrease constitutes the arrow of time.
    It does not require a separate thermodynamic argument or a special initial
    condition external to the framework; it is built into the admissibility structure of the code layer.
    The question ``why does time have a direction?'' thus reduces to
    ``why is the initial state~$U_0$ not already a global minimum of~$E_\chi$?''---a
    question about the boundary conditions of the cascade, not about the nature of time itself.

  \paragraph{The present as a dynamical frontier.}
    Within the cascade, the universe state~$U_n$ partitions the space of relational
    configurations into two complementary sectors:
    \begin{itemize}
      \item configurations already stabilised under~$\Pi$, whose projected images
      are fixed and no longer carry residual relational gradient;
      \item configurations that are not yet admissible, in the sense that the
      current spectral structure of~$\mathcal{L}$ does not yet support a locally injective projection onto them.
    \end{itemize}
    The present corresponds to the dynamical frontier between these two sectors:
    the region where the projection is becoming stable but where the cascade has not yet terminated.
    Because the non-injective fibre~$\Pi^{-1}(U_n)$ retains finite multiplicity at
    this frontier, the present is not an infinitesimal instant but a zone of finite
    relational extent, consistent with the finite propagation speeds imposed by the bounded-flux constraint.
    This provides a structural account of the physical thickness of the present
    without invoking additional postulates.

  \paragraph{Co-emergence of space and time from a single operator.}
    The preceding paragraphs have described space and time as having logically distinct
    origins within the framework: spatial structure follows from the admissibility
    envelopes of~$\chi$ and the fibre topology of~$\Pi$
    (see~\cite{Beau2026a1} and Section~\ref{sec:spectral-admissibility}),
    while temporal succession follows from the ordered execution of the cascade.
    One might therefore ask why space and time remain as tightly coupled as special relativity requires.

    The answer lies in the structure of the effective operator~$\mathcal{L}$.
    As established in~\cite[Sec.~6]{Beau2026c},
    the bounded-flux condition~$c_\chi$ constrains the admissible forms of the
    effective Laplacian~$\mathcal{L}$ simultaneously in its spatial and its cascade (temporal) sectors.
    Specifically:
    \begin{itemize}
      \item A purely elliptic (spatial) operator is excluded because it would admit
      instantaneous propagation, violating~$c_\chi$.
      \item An ultra-hyperbolic operator (more than one time-like direction) is
      excluded because it leads to an ill-posed Cauchy problem, destabilising the cascade.
      \item The unique compatible class is a hyperbolic operator with exactly one
      time-like direction and three space-like directions, whose principal symbol
      defines the Minkowski metric~$\eta_{\mu\nu} = \mathrm{diag}(-1,+1,+1,+1)$.
    \end{itemize}
    The Lorentz group therefore emerges not from a separate postulate but as the
    symmetry group of the unique operator class that is simultaneously compatible with
    three-dimensional spatial admissibility and with the irreversible, bounded-flux cascade.
    Space and time are logically separable in their origins---admissibility selects the
    former, cascade ordering generates the latter---but they are dynamically inseparable
    in the effective description, because only one operator class satisfies both constraints at once.
    This is why the universal ratio of spatial to temporal extents is fixed at~$c$,
    and why no admissible state of the cascade can disentangle space from time at the macroscopic level.

  \paragraph{Metric rigidity and the general-relativistic space--time budget.}
    The signature argument above operates at the level of special relativity:
    it establishes the existence of a pseudo-Riemannian metric~$g_{\mu\nu}$ with
    signature~$(-+++)$ and guarantees Lorentz invariance of the flat effective geometry.
    It does not, however, explain why the space--time budget remains exact and
    self-consistent in the presence of matter and curvature, i.e.\ why dilating the
    temporal metric coefficient at one point is compensated by a corresponding
    contraction of the spatial coefficients so that~$ds^2 = g_{\mu\nu}dx^\mu dx^\nu$
    remains a well-defined invariant throughout a curved spacetime.

    This stronger rigidity is enforced at a second level, through the spectral entropy
    functional
    \begin{equation}
      S_\Pi[g] \;=\; \tfrac{1}{2}\log\det' A_g,
    \end{equation}
    where~$A_g = -\nabla_g^2$ is the Laplace-type operator on the effective geometry
    (~\cite[Sec.~2, Sec.~4]{Beau2026h}).
    The renormalised metric variation of~$S_\Pi[g]$ generates, as its dominant infrared
    term, the Einstein tensor:
    \begin{equation}
      \frac{\delta S_\Pi^{\mathrm{ren}}}{\delta g^{\mu\nu}}
      \;=\; c_{\mathrm{EH}}\,G_{\mu\nu}
      + c_\Lambda\,g_{\mu\nu}
      + \text{(higher-derivative and non-local terms)}.
    \end{equation}
    The stationarity condition~$\delta S_\Pi[g]/\delta g^{\mu\nu} = 0$, imposed as the
    equilibrium requirement of the projection, therefore yields the Einstein field
    equations as the constraint that the effective metric must satisfy at each step of
    the cascade~\cite{Beau2026h, Beau2026i}.
    It is precisely these equations that lock the relationship between the temporal and
    spatial components of~$g_{\mu\nu}$ at every spacetime point: they prevent the metric
    from being deformed in a way that would violate the invariance of~$ds^2$.
    The~$A(r) = B(r)^{-1}$ identity of the Schwarzschild solution, which is the most
    explicit expression of this space--time budget compensation, is a direct consequence
    of flux conservation under the constrained effective operator~(Section~\cite[Sec.~7]{Beau2026c}
    and~\cite[Sec.~5.3, 6]{Beau2026h}).

    The logical chain is therefore the following.
    The bounded-flux condition~$c_\chi$ selects the $(-+++)$ signature, guaranteeing
    the existence of a pseudo-Riemannian metric and the universality of~$c$~(special relativity level).
    The spectral entropy functional~$S_\Pi[g]$, whose extremisation encodes the
    equilibrium of the cascade projection, then dynamically locks the metric coefficients
    through the Einstein equations, enforcing the exact compensation of the space--time
    budget under arbitrary matter distributions and curvature~(general relativity level).
    The two levels are not independent: both follow from the same bounded-projection
    structure of~$\chi$, operating at different orders in the effective gradient expansion.

  \paragraph{Relation to the block-universe picture.}
    The present account differs from a block-universe interpretation.
    A block universe asserts that all states~$U_n$ co-exist in a four-dimensional
    manifold, and that temporal flow is an illusion of perspective.
    In the present framework, the structural rules co-exist (the code is fixed),
    but the realisation of~$U_n$ is sequential and self-generated.
    No co-existence of all states is required or implied.
    Each state engenders its successor through its own admissibility structure.
    The appropriate analogy is a causal automaton rather than a static manifold:
    the program runs, step by step, driven by its own output.

  \paragraph{Summary.}
    The five elements that together constitute a logically closed account of temporal
    emergence and space--time coupling are:
    \begin{enumerate}
      \item \emph{$\chi$ as static code.}
      The admissibility rules, the Laplacian structure, and the saturation
      bound~$c_\chi$ are fixed and atemporal; they form the structural layer.
      \item \emph{$U_n$ as cursor on fibre classes.}
      The current universe state is self-referential: it determines, via its
      own spectral structure, which fibre classes are admissible successors.
      No external parameter advances the system.
      \item \emph{The bounded-flux constraint as the unique selector.}
      Among candidate transitions, only the steepest admissible descent along~$E_\chi$
      satisfies~$c_\chi$; the constraint itself acts as the selection mechanism,
      with no additional rule required.
      \item \emph{Co-emergence of space and time (special relativity level).}
      The spatial and temporal sectors are logically distinct in origin but
      dynamically inseparable: the bounded-flux operator uniquely selects the
      $(-+++)$ signature, fixing the universal ratio~$c$ and coupling the two
      sectors into Lorentz-invariant structure.
      \item \emph{Metric rigidity (general relativity level).}
      The spectral entropy functional~$S_\Pi[g]$, whose extremisation encodes the
      equilibrium of the cascade projection, generates the Einstein equations as its
      dominant infrared condition.
      These equations dynamically lock the relationship between spatial and temporal
      metric coefficients at every point, enforcing the exact compensation of the
      space--time budget---the invariance of~$ds^2$---under arbitrary matter
      distributions and curvature.
    \end{enumerate}
    \noindent
    Items 4 and~5 answer the question of space--time coupling at two distinct levels:
    item~4 explains why there is a single universal speed~$c$ and Lorentz invariance;
    item~5 explains why the metric compensation remains exact and position-dependent
    in the presence of gravity.
    Both levels follow from the same bounded-projection structure of~$\chi$, operating
    at different orders in the effective gradient expansion.
    Together, all five elements provide a logically closed account of temporal emergence
    that is consistent with the spatial structure established in the spectral admissibility
    programme~\cite{Beau2026a1}, with the Lorentz and Einstein structures
    derived in~\cite{Beau2026c, Beau2026h}, and with the relational fluctuation structure
    that underlies quantum indeterminacy in the framework.
